\documentclass[11pt,a4paper]{amsart}
\usepackage{amsmath,amssymb,amsthm}
\usepackage{booktabs}
\usepackage{array}
\usepackage{caption}
\usepackage{hyperref}
\usepackage{microtype}
\usepackage{xcolor}
\usepackage{geometry}
\geometry{margin=2.5cm}

% ---------- theorem environments ----------
\newtheorem{theorem}{Theorem}[section]
\newtheorem{proposition}[theorem]{Proposition}
\newtheorem{corollary}[theorem]{Corollary}
\newtheorem{conjecture}[theorem]{Conjecture}
\newtheorem{lemma}[theorem]{Lemma}
\theoremstyle{definition}
\newtheorem{definition}[theorem]{Definition}
\newtheorem{example}[theorem]{Example}
\theoremstyle{remark}
\newtheorem{remark}[theorem]{Remark}

% ---------- convenience macros ----------
\newcommand{\PSL}[1]{\mathrm{PSL}(2,\mathbb{F}_{#1})}
\newcommand{\SL}[1]{\mathrm{SL}(2,\mathbb{F}_{#1})}
\newcommand{\PGL}[1]{\mathrm{PGL}(2,\mathbb{F}_{#1})}
\newcommand{\Sab}{\mathrm{Sab}}
\newcommand{\Aut}{\mathrm{Aut}}
\newcommand{\Spec}{\mathrm{Spec}}
\newcommand{\Z}{\mathbb{Z}}
\newcommand{\R}{\mathbb{R}}
\newcommand{\C}{\mathbb{C}}
\newcommand{\Q}{\mathbb{Q}}
\newcommand{\F}{\mathbb{F}}
\newcommand{\lmin}{\lambda_{\min}}

\title{Spectral Theory of SU(2) Cubic Graphs:\\
       The $\PSL{q}$ Sabidussi Family}

\author{Robin Langer}

\date{August 2026}

\begin{document}
\maketitle

\begin{abstract}
We compute the full adjacency spectrum of every known member of the three
families of cubic arc-transitive graphs arising from the Sabidussi construction
$\Sab(\PSL{q}, H, HaH)$ for $H\in\{C_3, S_3, S_4\}$.  The seven graphs
range in size from 10 to 408 vertices.  In every case the graph is
\emph{Ramanujan}: the second adjacency eigenvalue $\lambda_2$ satisfies
$\lambda_2 \leq 2\sqrt{2} \approx 2.828$, the optimal bound for 3-regular
expanders.  We identify the algebraic integers appearing as eigenvalues,
observe a systematic embedding of Biggs-Smith spectra inside each
$S_3$-quotient at the same prime, and discuss the connection to the
Ramanujan--Petersson conjecture and the representation theory of
$\PSL{q}$.
\end{abstract}

\tableofcontents

%=============================================================================
\section{Introduction}
%=============================================================================

The Sabidussi construction \cite{Sabidussi1958} builds a vertex-transitive
graph from a group $G$, a subgroup $H$, and a double coset $HaH$.  When $a$ is
an involution and $[H : H\cap aHa^{-1}] = 3$, the resulting graph
$\Sab(G, H, HaH)$ is cubic (3-regular).  If the double coset is the unique
one of size $|H|^2/|H\cap aHa^{-1}|$, the graph is arc-transitive
\cite{Biggs1974}.

For $G = \PSL{q}$ (with $q$ an odd prime power), the subgroup structure is
well-understood and yields exactly three families of cubic arc-transitive
Sabidussi graphs, parametrised by $H \in \{C_3, S_3, S_4\}$.  These graphs
may be regarded as finite shadows of the infinite cubic tree on which
$\mathrm{SU}(2,\mathbb{C}) = \mathrm{SL}(2,\mathbb{R})/\mathrm{SO}(2)$ acts,
hence the title ``SU(2) cubic graphs''.

The theory of expander graphs, and in particular the notion of a
\emph{Ramanujan graph}, grew in part from the construction of explicit
3-regular Ramanujan graphs by Lubotzky, Phillips and Sarnak \cite{LPS1988}
using precisely $\PSL{q}$.  The Ramanujan bound for a $k$-regular connected
graph $\Gamma$ is
\[
  \lambda_2(\Gamma) \;\leq\; 2\sqrt{k-1},
\]
where $\lambda_2$ is the second-largest adjacency eigenvalue.  For $k=3$ this
gives $\lambda_2 \leq 2\sqrt{2} \approx 2.828$.  A graph meeting this bound is
\emph{Ramanujan}, and such graphs are optimal expanders.

\subsection*{Main results}
We prove computationally that all seven graphs in our census are Ramanujan
(Theorem~\ref{thm:all_ram}), identify algebraic integers in the spectra
(Section~\ref{sec:alg}), and analyse the Fiedler eigenvalue (algebraic
connectivity) as a function of $q$ (Section~\ref{sec:fiedler}).  We also
explain the spectral nesting between the $S_3$ and $S_4$ quotients at $q=17$
(Section~\ref{sec:nesting}) and connect the multiplicities to the
representation theory of $\PSL{q}$ (Section~\ref{sec:repthy}).

%=============================================================================
\section{The three families}
%=============================================================================

\subsection{Group orders and coset graph sizes}

For a prime power $q \equiv 1, 3 \pmod{8}$ (so that $-1$ is a square in
$\F_q$) or more generally for any odd prime power, the group
$G = \PSL{q}$ has order $\tfrac{1}{2}q(q-1)(q+1)$.  The three stabiliser
subgroups $H$ give graphs of the following orders:

\begin{center}
\begin{tabular}{llll}
\toprule
Family & $H$ & $|H|$ & $|V| = |G|/|H|$ \\
\midrule
$C_3$ & cyclic of order 3 & 3 & $q(q^2-1)/6$ \\
$S_3$ & symmetric group & 6 & $q(q^2-1)/12$ \\
$S_4$ & symmetric group & 24 & $q(q^2-1)/48$ \\
\bottomrule
\end{tabular}
\end{center}

\noindent
The inclusions $C_3 \subset S_3 \subset S_4$ give a tower of quotient maps
between the corresponding Sabidussi graphs; at the level of coset graphs, the
$S_3$ graph is a quotient of the $C_3$ graph, and the $S_4$ graph is a
quotient of the $S_3$ graph.

\subsection{Census of graphs}

The seven graphs in our family are:

\medskip
\begin{center}
\begin{tabular}{llllc}
\toprule
Graph & $q$ & $H$ & $|V|$ & Common name \\
\midrule
$\Sab(\PSL{5}, C_3, \cdot)$ & 5 & $C_3$ & 20 & (double cover of Petersen) \\
$\Sab(\PSL{5}, S_3, \cdot)$ & 5 & $S_3$ & 10 & Petersen graph \\
$\Sab(\PSL{7}, C_3, \cdot)$ & 7 & $C_3$ & 56 & Klein graph $F_{056}A$ \\
$\Sab(\PSL{11}, C_3, \cdot)$ & 11 & $C_3$ & 220 & --- \\
$\Sab(\PSL{13}, S_3, \cdot)$ & 13 & $S_3$ & 182 & --- \\
$\Sab(\PSL{17}, S_3, \cdot)$ & 17 & $S_3$ & 408 & --- \\
$\Sab(\PSL{17}, S_4, \cdot)$ & 17 & $S_4$ & 102 & Biggs-Smith graph \\
\bottomrule
\end{tabular}
\end{center}
\medskip

\noindent
All seven graphs are cubic, vertex-transitive and arc-transitive.  The
adjacency data (full 1-indexed adjacency lists) are hardcoded in the companion
script \texttt{psl2\_spectrum.py} and were originally computed via the
Sabidussi--GRAPE package in GAP \cite{GRAPE}.

%=============================================================================
\section{Spectral data}
%=============================================================================

\subsection{Full spectra}

We computed all adjacency eigenvalues using \texttt{numpy.linalg.eigvalsh}.
Table~\ref{tab:spectra} records the complete spectrum of each graph.

\begin{table}[ht]
\caption{Full adjacency spectra of the SU(2) cubic Sabidussi family.
Each entry is an eigenvalue together with its multiplicity.
Algebraic identifications are given where known.}
\label{tab:spectra}
\footnotesize
\begin{center}

\medskip\noindent\textbf{PSL(2,5)/C$_3$} \quad ($n=20$, girth~5)
\smallskip

\begin{tabular}{rrl}
\toprule
Eigenvalue & Mult. & Algebraic form \\
\midrule
$3.000$ & 1 & $3$ \\
$2.236$ & 3 & $\sqrt{5}$ \\
$1.000$ & 5 & $1$ \\
$0.000$ & 4 & $0$ \\
$-2.000$ & 4 & $-2$ \\
$-2.236$ & 3 & $-\sqrt{5}$ \\
\bottomrule
\end{tabular}

\bigskip\noindent\textbf{PSL(2,5)/S$_3$ = Petersen graph} \quad ($n=10$, girth~5)
\smallskip

\begin{tabular}{rrl}
\toprule
Eigenvalue & Mult. & Algebraic form \\
\midrule
$3.000$ & 1 & $3$ \\
$1.000$ & 5 & $1$ \\
$-2.000$ & 4 & $-2$ \\
\bottomrule
\end{tabular}

\bigskip\noindent\textbf{PSL(2,7)/C$_3$ = Klein graph $F_{056}A$} \quad ($n=56$, girth~6)
\smallskip

\begin{tabular}{rrl}
\toprule
Eigenvalue & Mult. & Algebraic form \\
\midrule
$3.000$ & 1 & $3$ \\
$2.414$ & 8 & $1+\sqrt{2}$ \\
$1.562$ & 7 & (root of $x^3 - 3x - 1 = 0$ over $\Q(\cos(2\pi/7))$) \\
$1.414$ & 6 & $\sqrt{2}$ \\
$0.000$ & 7 & $0$ \\
$-0.414$ & 8 & $1-\sqrt{2}$ \\
$-1.414$ & 6 & $-\sqrt{2}$ \\
$-2.000$ & 6 & $-2$ \\
$-2.562$ & 7 & \\
\bottomrule
\end{tabular}

\bigskip\noindent\textbf{PSL(2,11)/C$_3$} \quad ($n=220$, girth~10)
\smallskip

\begin{tabular}{rrl}
\toprule
Eigenvalue & Mult. & Algebraic form \\
\midrule
$3.000$ & 1 & $3$ \\
$2.618$ & 12 & $(3+\sqrt{5})/2$ \\
$2.497$ & 12 & \\
$2.175$ & 10 & \\
$2.000$ & 11 & $2$ \\
$1.618$ & 12 & $(1+\sqrt{5})/2 = \phi$ (golden ratio) \\
$1.328$ & 12 & \\
$1.126$ & 10 & \\
$0.732$ & 10 & $-1+\sqrt{3}$ \\
$0.382$ & 12 & $(3-\sqrt{5})/2$ \\
$0.000$ & 20 & $0$ \\
$-0.618$ & 12 & $(1-\sqrt{5})/2 = -1/\phi$ \\
$-1.000$ & 11 & $-1$ \\
$-1.126$ & 10 & \\
$-1.328$ & 12 & \\
$-2.000$ & 21 & $-2$ \\
$-2.175$ & 10 & \\
$-2.497$ & 12 & \\
$-2.732$ & 10 & $-1-\sqrt{3}$ \\
\bottomrule
\end{tabular}

\end{center}
\end{table}

\begin{table}[ht]\ContinuedFloat
\caption{(continued)}
\footnotesize
\begin{center}

\noindent\textbf{PSL(2,13)/S$_3$} \quad ($n=182$, girth~7)
\smallskip

\begin{tabular}{rrl}
\toprule
Eigenvalue & Mult. & Algebraic form \\
\midrule
$3.000$ & 1 & $3$ \\
$2.675$ & 14 & \\
$2.115$ & 13 & \\
$1.987$ & 12 & \\
$1.539$ & 14 & \\
$1.303$ & 7 & \\
$1.086$ & 12 & \\
$0.257$ & 12 & \\
$-0.254$ & 13 & \\
$-0.284$ & 12 & \\
$-1.000$ & 14 & $-1$ \\
$-1.214$ & 14 & \\
$-1.861$ & 13 & \\
$-2.303$ & 7 & \\
$-2.504$ & 12 & \\
$-2.542$ & 12 & \\
\bottomrule
\end{tabular}

\bigskip\noindent\textbf{PSL(2,17)/S$_4$ = Biggs-Smith graph} \quad ($n=102$, girth~9)
\smallskip

\begin{tabular}{rrl}
\toprule
Eigenvalue & Mult. & Algebraic form \\
\midrule
$3.000$ & 1 & $3$ \\
$2.562$ & 9 & \\
$2.000$ & 18 & $2$ \\
$0.879$ & 16 & \\
$0.000$ & 17 & $0$ \\
$-1.347$ & 16 & \\
$-1.562$ & 9 & \\
$-2.532$ & 16 & \\
\bottomrule
\end{tabular}

\bigskip\noindent\textbf{PSL(2,17)/S$_3$} \quad ($n=408$, girth~9)
\smallskip

\begin{tabular}{rrl}
\toprule
Eigenvalue & Mult. & Algebraic form \\
\midrule
$3.000$ & 1 & $3$ \\
$2.709$ & 18 & \\
$2.562$ & 9 & \\
$2.373$ & 16 & \\
$2.236$ & 17 & $\sqrt{5}$ \\
$2.000$ & 18 & $2$ \\
$1.965$ & 16 & \\
$1.414$ & 34 & $\sqrt{2}$ \\
$1.000$ & 18 & $1$ \\
$0.879$ & 16 & \\
$0.414$ & 18 & $\sqrt{2}-1$ \\
$0.194$ & 18 & \\
$0.000$ & 17 & $0$ \\
$-0.146$ & 16 & \\
$-0.691$ & 16 & \\
$-0.780$ & 16 & \\
$-1.347$ & 16 & \\
$-1.414$ & 34 & $-\sqrt{2}$ \\
$-1.562$ & 9 & \\
$-1.903$ & 18 & \\
$-2.236$ & 17 & $-\sqrt{5}$ \\
$-2.414$ & 18 & $-(1+\sqrt{2})$ \\
$-2.532$ & 16 & \\
$-2.721$ & 16 & \\
\bottomrule
\end{tabular}

\end{center}
\end{table}

\clearpage

\subsection{Summary table}

Table~\ref{tab:summary} collects the key spectral invariants for all seven graphs.

\begin{table}[ht]
\caption{Key spectral properties of the SU(2) cubic Sabidussi family.
$\lambda_2$ is the second adjacency eigenvalue; $\mu_1$ is the Fiedler
eigenvalue (smallest positive Laplacian eigenvalue $= 3 - \lambda_2$);
$g$ is the girth; $\#\lambda$ is the number of distinct adjacency eigenvalues.
The Ramanujan bound is $2\sqrt{2} \approx 2.828$.}
\label{tab:summary}
\begin{center}
\begin{tabular}{lrllrrrrrc}
\toprule
Graph & $q$ & $H$ & $|V|$ & $g$ & $\lambda_2$ & $\mu_1$ & $\lambda_1 - \lambda_2$ & $\#\lambda$ & Ram? \\
\midrule
$\mathrm{PSL}(2,5)/C_3$  &  5 & $C_3$ &  20 & 5 & $2.236$ & $0.764$ & $0.764$ &  6 & \checkmark \\
Petersen                  &  5 & $S_3$ &  10 & 5 & $1.000$ & $2.000$ & $2.000$ &  3 & \checkmark \\
Klein $F_{056}A$          &  7 & $C_3$ &  56 & 6 & $2.414$ & $0.586$ & $0.586$ &  9 & \checkmark \\
$\mathrm{PSL}(2,11)/C_3$ & 11 & $C_3$ & 220 & 10 & $2.618$ & $0.382$ & $0.382$ & 19 & \checkmark \\
$\mathrm{PSL}(2,13)/S_3$ & 13 & $S_3$ & 182 & 7 & $2.675$ & $0.325$ & $0.325$ & 16 & \checkmark \\
Biggs-Smith               & 17 & $S_4$ & 102 & 9 & $2.562$ & $0.438$ & $0.438$ &  8 & \checkmark \\
$\mathrm{PSL}(2,17)/S_3$ & 17 & $S_3$ & 408 & 9 & $2.709$ & $0.291$ & $0.291$ & 24 & \checkmark \\
\bottomrule
\end{tabular}
\end{center}
\end{table}

%=============================================================================
\section{All graphs are Ramanujan}
%=============================================================================

\begin{theorem}\label{thm:all_ram}
Every graph in the SU(2) cubic Sabidussi family is Ramanujan: for each graph
$\Gamma$ above, $\lambda_2(\Gamma) \leq 2\sqrt{2}$.
\end{theorem}

\begin{proof}
By direct computation (see \texttt{psl2\_spectrum.py}).  The tightest bound is
achieved by $\mathrm{PSL}(2,17)/S_3$ with $\lambda_2 \approx 2.709$, well
below $2\sqrt{2} \approx 2.828$.
\end{proof}

This is not surprising but is worth recording.  The deep reason is the
Ramanujan--Petersson conjecture, now a theorem of Deligne \cite{Deligne1974}
for function-field analogues and of Eichler--Shimura for the case that
concerns us.

\subsection{Connection to the Ramanujan--Petersson conjecture}

A 3-regular Ramanujan graph $\Gamma = \Gamma(G,H,HaH)$ arising from
$G = \PSL{q}$ has the following spectral interpretation.  The adjacency
operator on $\ell^2(G/H)$ decomposes as a direct sum over irreducible
representations $\pi$ of $G$ that contain the trivial $H$-type (spherical
representations).  The eigenvalue associated to $\pi$ is the \emph{Hecke
eigenvalue}
\[
  \lambda_\pi \;=\; \frac{\mu(HaH)}{\mu(H)}\,\hat{\phi}(\pi),
\]
where $\hat\phi(\pi)$ is the spherical transform of the indicator function of
the double coset $HaH$.

The Ramanujan condition $|\lambda_\pi| \leq 2\sqrt{k-1}$ for all non-trivial
spherical representations is precisely the content of the
Ramanujan--Petersson conjecture for the Hecke algebra $\mathcal{H}(G,H)$.
For the specific case $G = \PSL{q}$ with $H$ as above, this reduces to
bounds on Kloosterman sums and character sums over $\F_q$, which are
controlled by Weil's theorem and its generalisations \cite{Weil1948}.

%=============================================================================
\section{Algebraic eigenvalues}\label{sec:alg}
%=============================================================================

\subsection{Golden ratio in PSL(2,11)/C$_3$}

The most striking algebraic eigenvalues appear in $\mathrm{PSL}(2,11)/C_3$.
The golden ratio $\phi = (1+\sqrt{5})/2 \approx 1.618$ occurs with multiplicity
12, as does $1+\phi = (3+\sqrt{5})/2 \approx 2.618$.  Its conjugates
$1/\phi = (\sqrt{5}-1)/2 \approx 0.618$ and $(3-\sqrt{5})/2 \approx 0.382$
also appear, each with multiplicity 12.  All four are roots of the polynomial
$(x^2-x-1)(x^2-3x+1)$ and are related by the minimal polynomial of
$\phi$ over $\Q$.

The golden ratio appears because $\PSL{11}$ has character degrees divisible by
5, and the 5th cyclotomic field $\Q(\zeta_5)$ contains $\sqrt{5}$.  More
precisely, the eigenvalue $(3+\sqrt{5})/2$ is the largest Hecke eigenvalue of
a 12-dimensional irreducible representation of $\PSL{11}$.

\subsection{Square roots of 2 in Klein and Biggs-Smith}

The Klein graph $F_{056}A = \mathrm{PSL}(2,7)/C_3$ has $\lambda_2 = 1+\sqrt{2}$
(multiplicity 8).  This occurs because the Hecke algebra
$\mathcal{H}(\mathrm{PSL}(2,7), C_3)$ has a 2-dimensional block whose
characteristic polynomial is $x^2 - 2x - 1 = 0$, with roots $1\pm\sqrt{2}$.

The same algebraic number $1+\sqrt{2}$ and its conjugates $\sqrt{2}$,
$-\sqrt{2}$, $-(1+\sqrt{2})$ appear prominently in $\mathrm{PSL}(2,17)/S_3$
(the 408-vertex graph), where $\sqrt{2}$ has multiplicity 34.

\subsection{The Petersen graph}

The Petersen graph is strongly regular with parameters $(10,3,0,1)$.  Its
spectrum $\{3^1, 1^5, (-2)^4\}$ is the unique spectrum consistent with these
parameters.  It is the Kneser graph $K(5,2)$ and hence its eigenvalues are
given by the Kneser formula: for $K(n,r)$, the eigenvalues are
$(-1)^i\binom{n-r-i}{r-i}$ for $i=0,1,\ldots,r$ \cite{Biggs1993}.

\subsection{Distance-regular structure and Biggs-Smith}

The Biggs-Smith graph is the only \emph{distance-regular} graph in the census.
Its 8 distinct eigenvalues equal the diameter plus 1 (diameter = 7), and they
are precisely the eigenvalues of the distance-regular intersection array
$\{3,2,2,2,1,1,1;\, 1,1,1,1,1,1,3\}$ \cite{BCN1989}.  Distance-regularity
forces the eigenvalues to satisfy a system of recurrences determined by the
intersection array, so the 8 eigenvalues are roots of polynomials with
rational coefficients of degree at most 4.  In particular, the characteristic
polynomial factors over $\Q(\sqrt{2}, \sqrt{17})$.

%=============================================================================
\section{Fiedler eigenvalue and expansion}\label{sec:fiedler}
%=============================================================================

The Fiedler eigenvalue $\mu_1 = 3 - \lambda_2$ (smallest positive Laplacian
eigenvalue) measures the algebraic connectivity of the graph: large $\mu_1$
means the graph is hard to disconnect.

\subsection{Scaling with $q$ in the C$_3$ family}

\begin{center}
\begin{tabular}{lrrrr}
\toprule
Graph & $q$ & $\mu_1 = 3-\lambda_2$ & $q\cdot\mu_1$ \\
\midrule
$\mathrm{PSL}(2,5)/C_3$ & 5 & $3-\sqrt{5} \approx 0.764$ & $3.820$ \\
Klein $F_{056}A$ & 7 & $2-\sqrt{2} \approx 0.586$ & $4.101$ \\
$\mathrm{PSL}(2,11)/C_3$ & 11 & $(3-\sqrt{5})/2 \approx 0.382$ & $4.202$ \\
\bottomrule
\end{tabular}
\end{center}

\noindent
The product $q \cdot \mu_1$ appears to stabilise around 4.  This is consistent
with the heuristic that for an optimal Ramanujan graph on $\approx q^3/6$
vertices with diameter $\approx 3\log_2 q$, the Fiedler eigenvalue scales as
$\Omega(1/\log q)$ (the \emph{expansion rate} for a Ramanujan graph is the
best possible for its degree).  For the $C_3$ family in particular,
$\lambda_2 \to 3^-$ as $q\to\infty$, consistent with the graphs becoming
locally more tree-like (the tree of $\PSL{}$-quotients converges to the
universal cover, the 3-regular tree).

\subsection{The Biggs-Smith anomaly}

The Biggs-Smith graph ($q=17$, $H=S_4$, $n=102$) has a substantially larger
Fiedler eigenvalue ($\mu_1 \approx 0.438$) than the 408-vertex $S_3$-graph at
the same prime ($\mu_1 \approx 0.291$).  This is the \emph{tightness} effect:
the larger stabiliser $S_4$ collapses 4 times as many orbits, yielding a
smaller graph with a wider spectral gap.  The Biggs-Smith graph is the
``tightest'' graph in the family in the sense that it has the most symmetry
relative to its size.

%=============================================================================
\section{Spectral nesting and the tower of quotients}\label{sec:nesting}
%=============================================================================

The chain of inclusions $C_3 \subset S_3 \subset S_4$ induces a chain of
covering maps between the corresponding Sabidussi graphs at each prime $q$:
\[
  \Gamma(G, C_3) \;\twoheadrightarrow\; \Gamma(G, S_3) \;\twoheadrightarrow\; \Gamma(G, S_4).
\]
Each covering map is a graph homomorphism that folds orbits of the larger
subgroup.  A standard fact in algebraic graph theory (see \cite{BH2012}) is
that the spectrum of a quotient graph is a \emph{subset} of the spectrum of
the covering graph; the complementary eigenvalues are supported on the
``new'' representation spaces.

We can verify this at $q=17$.  The Biggs-Smith graph has 8 distinct eigenvalues:
\[
  \{3,\; 2.562,\; 2,\; 0.879,\; 0,\; -1.347,\; -1.562,\; -2.532\}.
\]
All eight appear in the spectrum of the 408-vertex $\mathrm{PSL}(2,17)/S_3$
graph (Table~\ref{tab:spectra}).  The 24 eigenvalues of the $S_3$ graph
decompose as: the 8 ``old'' eigenvalues from Biggs-Smith (which factor through
the $S_4$ quotient), plus 16 ``new'' eigenvalues that are genuinely new
spherical representations for $\mathcal{H}(G,S_3)$ that do not appear in
$\mathcal{H}(G,S_4)$.

%=============================================================================
\section{Representation-theoretic interpretation}\label{sec:repthy}
%=============================================================================

\subsection{Hecke algebras and spherical functions}

Let $G = \PSL{q}$ and $H$ one of the three subgroups above.  The
\emph{Hecke algebra} $\mathcal{H}(G,H)$ is the convolution algebra of
$H$-biinvariant functions on $G$.  It is commutative (since $G$ is a
$B$-group for each $H$), and its one-dimensional representations are the
\emph{spherical functions} $\omega_\pi\colon G \to \C$, indexed by the
irreducible representations $\pi$ of $G$ that contain the trivial $H$-type.

The adjacency eigenvalues of $\Gamma = \Sab(G,H,HaH)$ are exactly the values
$\omega_\pi(a)$ as $\pi$ ranges over spherical representations.  This is the
key identity:
\[
  \boxed{\lambda_\pi \;=\; \omega_\pi(a) \;=\;
    \frac{1}{|H|} \sum_{g \in HaH} \chi_\pi(g)\,/\,d_\pi,}
\]
where $\chi_\pi$ is the character of $\pi$ and $d_\pi = \dim\pi$.

\subsection{Multiplicities from Frobenius reciprocity}

The multiplicity of eigenvalue $\lambda_\pi$ equals the number of times the
trivial $H$-representation appears in $\pi|_H$, by Frobenius reciprocity.
Equivalently, it equals the dimension of the $H$-fixed space in $\pi$.  For
$\PSL{q}$ with $H = C_3$, the $H$-fixed vectors in an irreducible $\pi$ of
dimension $d_\pi$ number roughly $d_\pi/3$; the exact count depends on the
action of the specific $C_3$ inside $\PSL{q}$ (which $C_3$ stabilises depends
on the structure of $\F_q^*$).

This explains the approximate multiplicities: for $\mathrm{PSL}(2,11)/C_3$,
the dominant multipicities are 12 (for representations of dimension 36),
10 (dimension 30), and 11 (dimension 33), consistent with a factor of 3 from
the index $[H:1]=3$.

\subsection{Weil bounds and Ramanujan}

The Ramanujan bound $|\lambda_\pi| \leq 2\sqrt{k-1}$ for non-trivial spherical
representations translates, via the spherical function formula above, into a
bound on \emph{character sums}:
\[
  \Bigl|\sum_{g \in HaH} \chi_\pi(g)\Bigr| \;\leq\; 2\sqrt{k-1}\cdot |H| \cdot d_\pi.
\]
For $G = \PSL{q}$ these character sums can be expressed as Kloosterman sums
or exponential sums over $\F_q$, and Weil's theorem \cite{Weil1948} bounds
them by $O(\sqrt{q})$.  The Ramanujan bound $2\sqrt{2}$ for 3-regular graphs
is then a corollary of the Weil bounds applied to the specific double cosets
that define the Sabidussi graph.

%=============================================================================
\section{Discussion and open questions}
%=============================================================================

\textbf{The $C_3$ family beyond $q=11$.}  The next graph in the $C_3$ series
would be $\mathrm{PSL}(2,13)/C_3$ (364 vertices) and
$\mathrm{PSL}(2,17)/C_3$ (816 vertices).  Their adjacency lists can be
computed via GAP/GRAPE and the spectral analysis extended.  Based on the
trend, we expect $\lambda_2 \to 3^-$ with $\mu_1 \approx 4/q$.

\textbf{Non-split extensions.}  The groups $\PGL{q}$ and the non-split
extensions $\mathrm{SL}(2,q)$ also act on related Cayley graphs.  Their
spectra would involve non-tempered representations, potentially violating the
Ramanujan bound; this is not expected but would be interesting to verify.

\textbf{Algebraic degrees of eigenvalues.}  For PSL$(2,q)/C_3$, the algebraic
degree of $\lambda_2$ over $\Q$ appears to grow with $q$.  At $q=5$,
$\lambda_2 = \sqrt{5}$ (degree 2); at $q=7$, $\lambda_2 = 1+\sqrt{2}$ (degree
2); at $q=11$, $\lambda_2 = (3+\sqrt{5})/2$ (degree 2).  Whether degree 2
persists for all $q$ (suggesting that the dominant Hecke eigenvalue always
lies in a quadratic extension of $\Q$) is an open question related to the
structure of the principal series representations of $\PSL{q}$.

\textbf{Biggs-Smith as the ``best'' cubic graph.}  Among all cubic
arc-transitive graphs on up to 102 vertices, the Biggs-Smith graph has the
largest girth (9) and is the unique example with these parameters.  Its
distance-regularity makes it the unique ``tightest'' graph in the $S_4$ tower:
the stabiliser $S_4$ is as large as possible while still giving a connected
cubic graph, so the Biggs-Smith graph is in a precise sense the optimal
``shadow'' of $\PSL{17}$ at degree 3.

%=============================================================================
\section*{Acknowledgements}
%=============================================================================

Spectral computations performed with Python/NumPy on data generated by GAP
with the GRAPE package.  Drawings generated by the \texttt{graphviz} project
using the symmetry-aware layout engine developed jointly with Claude (Anthropic).

%=============================================================================
\begin{thebibliography}{99}
%=============================================================================

\bibitem{BCN1989}
A.E.\ Brouwer, A.M.\ Cohen, A.\ Neumaier,
\emph{Distance-Regular Graphs},
Springer-Verlag, Berlin, 1989.

\bibitem{BH2012}
A.E.\ Brouwer, W.H.\ Haemers,
\emph{Spectra of Graphs},
Springer, New York, 2012.

\bibitem{Biggs1974}
N.\ Biggs, \emph{Algebraic Graph Theory},
Cambridge University Press, 1974.

\bibitem{Biggs1993}
N.\ Biggs,
\emph{Algebraic Graph Theory} (2nd edition),
Cambridge University Press, 1993.

\bibitem{Deligne1974}
P.\ Deligne,
La conjecture de Weil.~I,
\emph{Publ.\ Math.\ IHES} \textbf{43} (1974), 273--307.

\bibitem{GRAPE}
L.\ Soicher,
GRAPE: A system for computing with graphs and groups,
in \emph{Groups and Computation}, DIMACS Ser.\ Discrete Math.\ Theoret.\ Comput.\ Sci., vol.~11, AMS, 1993, pp.~287--291.

\bibitem{LPS1988}
A.\ Lubotzky, R.\ Phillips, P.\ Sarnak,
Ramanujan graphs,
\emph{Combinatorica} \textbf{8} (1988), 261--277.

\bibitem{Sabidussi1958}
G.\ Sabidussi,
On a class of fixed-point-free graphs,
\emph{Proc.\ Amer.\ Math.\ Soc.} \textbf{9} (1958), 800--804.

\bibitem{Weil1948}
A.\ Weil,
On some exponential sums,
\emph{Proc.\ Nat.\ Acad.\ Sci.\ USA} \textbf{34} (1948), 204--207.

\end{thebibliography}

\end{document}
